\documentclass[dvipsnames]{book}

\usepackage[margin=1cm,a5paper]{geometry}

\usepackage{bold-extra}
\usepackage{polski}
\usepackage{iwona}
\usepackage[OT4]{fontenc}

\usepackage{amsmath,amssymb,amsthm}
\usepackage{pgf,tikz}
\usepackage{hyperref}

\pagestyle{empty}
\setlength{\parindent}{0pt}
\setlength{\parskip}{5pt}
\renewcommand\leq\leqslant
\renewcommand\geq\geqslant

\colorlet{maincolor}{ForestGreen}
\usepackage[most]{tcolorbox}
\usepackage{varwidth}

\newtcolorbox{tresc}[3][]{enhanced, sharp corners, boxrule=0mm, colback=maincolor!10,
colframe=black,toptitle=1mm,fonttitle=\bfseries,
colbacktitle=maincolor!25, coltitle=black,
title={Zadanie #2\hfill{\it Opracowanie: #3}},#1}

\newenvironment{problem}[3]{\newpage\begin{tresc}{#1}{#2}{#3}}{\end{tresc}}

% Dodatkowe makra

\usepackage{multicol}
\usepackage{cancel}
\usepackage{bbm}

\def\ind{\mathbbm{1}}

\def\N{\mathbb{N}}
\def\P{\mathbb{P}}
\def\E{\mathbb{E}}
\def\Z{\mathbb{Z}}
\def\Q{\mathbb{Q}}
\def\R{\mathbb{R}}

\def\nwd{\mathrm{NWD}}

\newtheorem*{lemat}{Lemat}

\usepackage{enumitem}
\setlist[enumerate,1]{label={(\alph*)},topsep=0pt}

% ------------------------------------------------------------------
% Kolorowe boxy: definicje (zielone), twierdzenia/obserwacje/wnioski (niebieskie),
% algorytmy (liliowe)
% ------------------------------------------------------------------

\colorlet{defcolor}{ForestGreen}
\colorlet{thmcolor}{RoyalBlue}
\colorlet{algcolor}{Orchid}

\newcounter{defcnt}
\newcounter{thmcnt}
\newcounter{algcnt}

% wspólny styl boxa
\tcbset{notatkibox/.style={enhanced, sharp corners, boxrule=0mm, frame hidden,
  toptitle=0.5mm, bottomtitle=0.5mm, left=2mm, right=2mm, top=1.5mm, bottom=1.5mm,
  fonttitle=\bfseries, coltitle=black, breakable, before skip=6pt, after skip=6pt}}

% \begin{definicja}[nazwa]
\NewDocumentEnvironment{definicja}{o}{%
  \refstepcounter{defcnt}%
  \begin{tcolorbox}[notatkibox, colback=defcolor!8, colframe=defcolor!25,
    colbacktitle=defcolor!25,
    title={Definicja \thedefcnt\IfValueT{#1}{\ (#1)}}]%
}{\end{tcolorbox}}

% generator środowisk niebieskich: #1 = nazwa środowiska, #2 = etykieta
\newcommand{\newthmbox}[2]{%
  \NewDocumentEnvironment{#1}{o}{%
    \refstepcounter{thmcnt}%
    \begin{tcolorbox}[notatkibox, colback=thmcolor!8, colframe=thmcolor!25,
      colbacktitle=thmcolor!25,
      title={#2 \thethmcnt\IfValueT{##1}{\ (##1)}}]%
  }{\end{tcolorbox}}%
}

\newthmbox{twierdzenie}{Twierdzenie}
\newthmbox{obserwacja}{Obserwacja}
\newthmbox{wniosek}{Wniosek}
\newthmbox{stwierdzenie}{Stwierdzenie}

% \begin{algorytm}[nazwa] -- liliowy box
\NewDocumentEnvironment{algorytm}{o}{%
  \refstepcounter{algcnt}%
  \begin{tcolorbox}[notatkibox, colback=algcolor!8, colframe=algcolor!25,
    colbacktitle=algcolor!25,
    title={Algorytm \thealgcnt\IfValueT{#1}{\ (#1)}}]%
}{\end{tcolorbox}}

% treść zadania -- żółty box, \begin{zadanie}{1.1}
\colorlet{zadcolor}{Dandelion}

\NewDocumentEnvironment{zadanie}{m}{%
  \begin{tcolorbox}[notatkibox, colback=zadcolor!15, colframe=zadcolor!40,
    colbacktitle=zadcolor!40,
    title={Zadanie #1}]%
}{\end{tcolorbox}}

% przykład -- bez boxa, wyróżniony nagłówkiem
\theoremstyle{definition}
\newtheorem*{przyklad}{Przykład}

\renewcommand{\qedsymbol}{$\square$}


\begin{document}

\begin{center}
  {\Large\bfseries Optymalizacja stochastyczna --- Wykład 1 --- 2026-10-05}
\end{center}

\section*{Sprawy organizacyjne}

\begin{itemize}[topsep=0pt,itemsep=0pt]
  \item Wykład: 9:30--11:00 \quad (sala 105)
  \item Ćwiczenia: 11:15--12:00 \quad (sala 105)
  \item Laboratoria: 12:30--14:00 \quad (sala 219)
\end{itemize}

Punktacja (min. 50 pkt.\ na ocenę 3.0):
\begin{itemize}[topsep=0pt,itemsep=0pt]
  \item Egzamin: 0--50 pkt.
  \item Kolokwium: 0--20 pkt.
  \item 5$\times$ zadanie z labów: 0--20 pkt.
  \item 4$\times$ kartkówki: 0--10 pkt.
  \item Aktywność na ćwiczeniach: 0--10 pkt.
\end{itemize}

Strona prowadzącego:
\url{http://pages.mini.pw.edu.pl/~joziakp/os}

\paragraph{Literatura:}
\begin{itemize}[topsep=0pt,itemsep=0pt]
  \item 1--5: różne źródła,
  \item 6--10: \emph{Introduction to Stochastic Search \& Optimization} --- James C. Spall,
  \item 11--15: \emph{Monte Carlo Statistical Methods} --- Ch.\ Robert, G.\ Casella.
\end{itemize}

\hrulefill

\section*{Czym zajmuje się optymalizacja stochastyczna?}

Zagadnienia postaci:
\[
  x^* = \operatorname*{arg\,min}_{x \in S} f(x).
\]

\begin{enumerate}[label=\arabic*)]
  \item $\nabla f(x) = 0$ $\longrightarrow$ zbiór punktów krytycznych
  $\longrightarrow$ wśród nich może być minimum.

  Ale: równanie $\nabla f(x) = 0$ może być bardzo trudne do rozwiązania
  analitycznie.

  \item Przybliżone rozwiązanie można uzyskać metodą kroku gradientowego.

  Ale:
  \begin{itemize}[topsep=0pt,itemsep=0pt]
    \item zbieżność tylko do minimum lokalnego,
    \item $\nabla$ może być trudny/niemożliwy do policzenia.
  \end{itemize}
\end{enumerate}

\begin{przyklad}
Szukamy $x^* = \operatorname*{arg\,min}_{x \in S} f(x)$, ale obserwujemy jedynie
$g(x, \omega)$ takie, że $f(x) = \E\big[g(x,\omega)\big]$, gdzie $\omega$ to
element przestrzeni probabilistycznej.

Np. $g(x,\omega) = f(x) + Z(\omega)$, gdzie $Z(\omega) \sim \mathcal{N}(0,\nu^2)$
jest szumem.
\end{przyklad}

\hrulefill

\begin{algorytm}[Przeszukiwanie Stochastyczne --- PRS, \emph{Pure Random Search}]
Znaleźć $x^* = \operatorname*{arg\,min}_{x \in S} f(x)$, gdzie $f$ jest ciągła,
$S \subseteq \R^d$.

Weźmy rozkład $\delta \colon \mathcal{B}(S) \to [0,1]$.

\begin{itemize}[topsep=0pt,itemsep=0pt]
  \item \textbf{Krok 0:} $y_0 \sim \delta$, \ $x_0 = y_0$.
  \item \textbf{Krok 1:} $y_{k+1} \sim \delta$, \quad
  $x_{k+1} = \begin{cases}
    y_{k+1} & : f(y_{k+1}) < f(x_k), \\
    x_k & : \text{wpp.}
  \end{cases}$
  \item \textbf{Krok 2:} $k \mathrel{{+}{=}} 1$.
\end{itemize}
Przerywamy, gdy $k$ przekroczy ustalony wcześniej limit.
\end{algorytm}

\section*{Model matematyczny PRS}

Niech $X_k = X_k(\omega)$, $Y_k = Y_k(\omega)$ dla pewnej $\omega \in \Omega$.

Zakładamy, że $Y_0, Y_1, \ldots \overset{\text{iid}}{\sim} \delta$ oraz
\[
  X_{k+1} = \begin{cases}
    Y_{k+1} & : f(Y_{k+1}) < f(X_k), \\
    X_k & : \text{wpp.}
  \end{cases}
\]

Oznaczamy:
\begin{gather*}
  m := \min_{x \in S} f(x), \\
  S_\varepsilon := \{x \in S : f(x) < m + \varepsilon\}, \\
  X^* := \{x \in S : f(x) = m\}.
\end{gather*}

\begin{twierdzenie}
Jeśli $f$ ciągła, to przy powyższych oznaczeniach:
\begin{gather*}
  \P\big(X_k \xrightarrow{k \to \infty} X^*\big) = 1, \\
  \P\big(\{\omega \in \Omega : d(X_k(\omega), X^*) \xrightarrow{k \to \infty} 0\}\big) = 1.
\end{gather*}
\end{twierdzenie}

\begin{proof}
\begin{enumerate}[label=\arabic*)]
  \item Ciąg $f(X_k)$ jest nierosnący i ograniczony z dołu przez $m$ (dla
  każdego $\omega \in \Omega$), więc z twierdzenia Bolzano--Weierstrassa jest
  zbieżny do $\eta(\omega) \geq m$.

  \item Ustalmy $\varepsilon > 0$. Niech
  \[
    R_k = \{\omega \in \Omega : f(Y_k(\omega)) < m + \varepsilon\} = Y_k^{-1}(S_\varepsilon)
  \]
  --- zdarzenia niezależne.

  \item $\P(R_k) = \P\big(Y_k^{-1}(S_\varepsilon)\big) = \delta(S_\varepsilon) =: \alpha > 0$.

  Z Lematu Borela--Cantelliego: $(\Omega, \mathcal{F}, \P)$ --- przestrzeń
  probabilistyczna, $R_1, R_2, \ldots$ --- zdarzenia niezależne,
  \[
    \sum_{k \geq 1} \P(R_k) = +\infty
    \implies
    \P\Big(\bigcap_{t=1}^{\infty} \bigcup_{i=t}^{\infty} R_i\Big) = 1.
  \]

  \item $\eta \geq m + \varepsilon \implies \forall k\ f(Y_k) \geq \eta \geq m + \varepsilon$, zatem
  \begin{gather*}
    \{\omega \in \Omega : \eta(\omega) \geq m + \varepsilon\}
      \subseteq \bigcap_{k=1}^{\infty} (\Omega \setminus R_k)
      = \Omega \setminus \bigcup_{k=1}^{\infty} R_k, \\
    \{\omega \in \Omega : \eta(\omega) \leq m + \varepsilon\}
      \supseteq \bigcup_{k=1}^{\infty} R_k
      \supseteq \bigcap_{t=1}^{\infty} \bigcup_{i=t}^{\infty} R_t,
  \end{gather*}
  więc $\P\big(\{\omega \in \Omega : \eta(\omega) \leq m + \varepsilon\}\big) = 1$.
\end{enumerate}
\end{proof}

\paragraph{Zalety PRS:}
\begin{itemize}[topsep=0pt,itemsep=0pt]
  \item łatwość implementacji,
  \item używa tylko pomiarów funkcji $f$,
  \item rozsądny koszt obliczeniowy,
  \item ogólność,
  \item mocne ugruntowanie teoretyczne.
\end{itemize}

\hrulefill

\begin{algorytm}[Localized Random Search]
\begin{itemize}[topsep=0pt,itemsep=0pt]
  \item \textbf{Krok 0:} ustalam punkt początkowy $x_0$ (losowo albo na podstawie
  wiedzy o problemie).
  \item \textbf{Krok 1:} losujemy $d_k \in \R^d$. Jeśli $x_k + d_k \in S$, to
  $y_{k+1} = x_k + d_k$.

  Wpp.\ $y_{k+1} = \operatorname*{arg\,min}_{x \in S} \|x - (x_k + d_k)\|$
  (najbliższy $x_k + d_k$ punkt z $S$).
  \item \textbf{Krok 2:} $x_{k+1} = \begin{cases}
    y_{k+1} & : f(y_{k+1}) < f(x_k), \\
    x_k & : \text{wpp.}
  \end{cases}$
  \item \textbf{Krok 3:} $k \mathrel{{+}{=}} 1$
\end{itemize}
Całość powtarzamy dopóki nie przekroczymy ustalonego limitu na $k$.
\end{algorytm}

Typowo $d_k \sim \mathcal{N}_d$ (wielowymiarowy normalny) oraz wariancja maleje ze
wzrostem $k$. Wtedy algorytm stabilizuje się, mamy mniejsze ,,skoki''.

\begin{twierdzenie}
$f \colon S \to \R$ ciągła, $m := \min_{x \in S} f(x)$,
$X^* := \{x \in S : f(x) = m\}$.

Rozważmy wariant algorytmu z $d_k \sim \mathcal{N}_d(0, I_d)$. Losujemy, dopóki
$x_k + d_k \in S$. Niech $\varepsilon > 0$, niech
$S_\varepsilon := \{x \in S : f(x) < m + \varepsilon\}$. Wtedy
\[
  \lim_{n \to \infty} \P\big(X_k \in S_\varepsilon\big) = 1.
\]
\end{twierdzenie}

\hrulefill

\begin{algorytm}[Poprawiony Localized Random Search]
\begin{itemize}[topsep=0pt,itemsep=0pt]
  \item \textbf{Krok 0:} ustal $x_0 \in S$ losowo bądź na podstawie wiedzy o
  problemie. $k = 0$, $b_k = 0$.
  \item \textbf{Krok 1:} losujemy niezależnie $d_k \sim \mathcal{N}_d(b_k, \sigma^2 I_d)$.
  Jeśli $x_k + d_k \notin S$, to losujemy ponownie lub rzutujemy $x_k + d_k$ na $S$.
  \item \textbf{Krok 2:} (wektor)
  \[
    \begin{bmatrix} x_{k+1} \\ b_{k+1} \end{bmatrix} =
    \begin{cases}
      \begin{bmatrix} x_k + d_k \\[1mm] (2d_k - b_k)/5 \end{bmatrix}
        & \text{jeśli } f(x_k + d_k) < f(x_k), \\[8mm]
      \begin{bmatrix} x_k + 2b_k - d_k \\[1mm] (7b_k - 2d_k)/5 \end{bmatrix}
        & : f(x_k + 2b_k - d_k) < f(x_k), \\[8mm]
      \begin{bmatrix} x_k \\[1mm] \tfrac{1}{2} b_k \end{bmatrix}
        & \text{wpp.}
    \end{cases}
  \]
  \item \textbf{Krok 3:} $k \mathrel{{+}{=}} 1$; jeśli kryterium stopu
  $\longrightarrow$ koniec, wpp.\ $\longrightarrow$ Krok 1.
\end{itemize}
\end{algorytm}

\hrulefill

\section*{Tempo zbieżności PRS I}

Naszym celem jest szacowanie liczby kroków, po której $X_k \in S_\varepsilon$
($\varepsilon > 0$ ustalony, reprezentuje tolerancję na odchylenie od minimum).
Oznaczamy
\[
  \rho(\varepsilon) := \delta(S_\varepsilon)
\]
i zakładamy, że $\rho(\varepsilon) > 0$ (w przeciwnym razie algorytm nigdy nie
trafi do $S_\varepsilon$).

\paragraph{Redukcja do prób Bernoulliego.}
Z definicji $X_k$ mamy $f(X_k) = \min_{j \leq k} f(Y_j)$ (indukcja po $k$: w
każdym kroku zostawiamy lepszy z punktów $Y_{k+1}$, $X_k$). Przynależność do
$S_\varepsilon$ zależy tylko od wartości $f$, więc
\[
  \{X_k \in S_\varepsilon\} = \bigcup_{j \leq k} \{Y_j \in S_\varepsilon\},
\]
czyli ,,do kroku $k$ trafiliśmy do $S_\varepsilon$'' $\iff$ ,,choć jedno
losowanie $Y_j$ wpadło do $S_\varepsilon$''. Zdarzenia
$\{Y_k \in S_\varepsilon\}$ są niezależne i każde ma prawdopodobieństwo
$\delta(S_\varepsilon) = \rho(\varepsilon)$, zatem szansa, że nie odwiedzimy
tego zbioru w krokach $1, 2, \ldots, n$, wynosi
\[
  \P\big(X_n \notin S_\varepsilon\big)
  = \prod_{k=1}^{n} \big(1 - \P(Y_k \in S_\varepsilon)\big)
  = \big(1-\rho(\varepsilon)\big)^n.
\]

\paragraph{Liczba kroków przy zadanej pewności.}
Ile kroków potrzeba, by trafić do $S_\varepsilon$ z prawdopodobieństwem
$\geq 1 - \gamma$ (dla małego $\gamma \in (0,1)$)? Ponieważ
$\log\big(1-\rho(\varepsilon)\big) < 0$, dzielenie przez tę liczbę odwraca
nierówność:
\begin{align*}
  1 - \big(1-\rho(\varepsilon)\big)^n \geq 1 - \gamma
  &\iff \big(1-\rho(\varepsilon)\big)^n \leq \gamma \\
  &\iff n \log\big(1-\rho(\varepsilon)\big) \leq \log \gamma \\
  &\iff n \geq \frac{\log \gamma}{\log\big(1-\rho(\varepsilon)\big)}.
\end{align*}
Oba logarytmy są ujemne, więc prawa strona jest dodatnia. Z nierówności
$\log(1-t) \leq -t$ dla $t \in [0,1)$ dostajemy wygodniejsze oszacowanie
\[
  \frac{\log \gamma}{\log\big(1-\rho(\varepsilon)\big)}
  \leq \frac{\log(1/\gamma)}{\rho(\varepsilon)},
\]
czyli koszt rośnie jak $1/\rho(\varepsilon)$, a od żądanej pewności zależy
jedynie logarytmicznie.

\paragraph{Czas pierwszego trafienia.}
Niech
\[
  N(\varepsilon) := \min\{k : X_k \in S_\varepsilon\}
  = \min\{k : Y_k \in S_\varepsilon\}
  = \min\{k : f(Y_k) < m + \varepsilon\}
\]
(druga równość to znów fakt, że $f(X_k)$ jest minimum dotychczasowych wartości).
Zatem $N(\varepsilon)$ to numer pierwszego sukcesu w ciągu niezależnych prób
Bernoulliego z prawdopodobieństwem sukcesu $\rho(\varepsilon)$:
\[
  \P\big(N(\varepsilon) = n\big) = \big(1-\rho(\varepsilon)\big)^{n-1} \rho(\varepsilon),
\]
czyli $N(\varepsilon) \sim \mathrm{Geo}\big(\rho(\varepsilon)\big)$, skąd
\[
  \E\big[N(\varepsilon)\big] = \frac{1}{\rho(\varepsilon)},
  \qquad
  \mathrm{Var}\big[N(\varepsilon)\big] = \frac{1 - \rho(\varepsilon)}{\rho(\varepsilon)^2}.
\]

\paragraph{Przypadek rozkładu jednostajnego.}
Niech $S \subseteq \R^d$ będzie ograniczony o mierze $|S| \in (0,\infty)$, a
$\delta$ --- rozkładem jednostajnym na $S$, tzn.
$\rho(\varepsilon) = |S_\varepsilon| / |S|$. Ustalmy $x^* \in X^*$; z ciągłości
$f$ istnieje $\varepsilon' > 0$ takie, że $B(x^*, \varepsilon') \subseteq S_\varepsilon$.
Ponieważ $|B(x^*, \varepsilon')| = c_d \cdot (\varepsilon')^d$, mamy
$\rho(\varepsilon) \geq c_d (\varepsilon')^d / |S|$, a więc
\[
  \E\big[N(\varepsilon)\big] \leq \frac{|S|}{c_d \cdot (\varepsilon')^d},
\]
z równością, gdy $S_\varepsilon = B(x^*, \varepsilon')$. Kluczowa jest tu
zależność $(\varepsilon')^{-d}$: oczekiwany czas rośnie wykładniczo z wymiarem
$d$ (,,przekleństwo wymiaru''), np.\ dziesięciokrotne zmniejszenie promienia
tolerancji mnoży koszt przez $10^d$.

\hrulefill

\begin{algorytm}[Pure Adaptive Search]
\begin{itemize}[topsep=0pt,itemsep=0pt]
  \item \textbf{Krok 0:} ustalmy $x_0 \in S$ (losowo lub na podstawie wiedzy o
  problemie). $k = 0$, $w_0 = f(x_0)$.
  \item \textbf{Krok 1:} losuj $x_{k+1}$ z rozkładu jednostajnego na
  \[
    S_k = \{x \in S : f(x) < w_k\}.
  \]
  \item \textbf{Krok 2:} $k \mathrel{{+}{=}} 1$. Jeśli kryterium stopu spełnione
  $\longrightarrow$ koniec, wpp.\ $\longrightarrow$ Krok 1.
\end{itemize}
\end{algorytm}

\hrulefill

\section*{Teoria rekordów}

Niech $X_1, X_2, \ldots$ --- ciąg zmiennych losowych iid o wspólnej dystrybuancie
$F(x)$.

Konstruujemy ciąg zmiennych losowych $N(k)$ w następujący sposób:
\begin{align*}
  N(1) &= 1, \\
  N(k) &= \min\big\{j > N(k-1) : X_j > X_{N(k-1)}\big\} \\
       &= \min\big\{j > N(k-1) : X_j > \max(X_1, \ldots, X_{N(k-1)})\big\}.
\end{align*}

Wtedy $\P\big(N(k) < \infty\big) = 1$, zatem
$\forall k\ \ p_{n,k} := \P\big(N(k) = n\big)$, więc $p_{n,k}$ jest rozkładem
prawdopodobieństwa na $\N$.

\begin{lemat}
Jeśli $F(x)$ to funkcja ciągła, to rozkład $N(k)$ nie zależy od $F$.
\end{lemat}

\begin{proof}
Ćwiczenie.
\end{proof}

\end{document}
